\section{The problem and its optimality conditions}
\label{sec:problem}

Write the program~\eqref{eq:qp} as
\begin{equation}
  \min_{x \in \R^n} \; f(x) := \tfrac12 x\T G x - a\T x
  \qquad \text{subject to} \qquad
  \begin{cases}
    c_i\T x = b_i, & i \le m_{\mathrm{eq}},\\
    c_i\T x \ge b_i, & i > m_{\mathrm{eq}},
  \end{cases}
  \label{eq:qp2}
\end{equation}
with $G \succ 0$. Two conventions are worth flagging because they differ from the
textbook ones and are inherited from the reference implementation's signature.
The linear term is \emph{subtracted}, so that the unconstrained minimiser is
$G^{-1}a$ rather than $-G^{-1}a$; and the constraints are held column-wise in
$C$, so that $C\T x \ge b$ rather than $Cx \le b$. Neither affects anything
below beyond signs.

\subsection{Existence, uniqueness, sufficiency}

Let $\Omega = \{x : c_i\T x = b_i \ (i \le m_{\mathrm{eq}}),\ c_i\T x \ge b_i \
(i > m_{\mathrm{eq}})\}$ denote the feasible set, a closed convex polyhedron.

Since $G \succ 0$ the objective is strictly convex and coercive, so \eqref{eq:qp2} has exactly one solution whenever $\Omega \neq \emptyset$; see~\cite[Ch.~16]{nocedal2006}. What the method needs beyond that is the converse of the usual necessary conditions, and that is worth proving because every certificate in this paper rests on it.

The Lagrangian of~\eqref{eq:qp2} is
$L(x,\lambda) = \tfrac12 x\T G x - a\T x - \lambda\T(C\T x - b)$, and the
Karush--Kuhn--Tucker conditions are
\begin{align}
  G x - a &= C \lambda, \label{eq:stat}\\
  c_i\T x &= b_i \quad (i \le m_{\mathrm{eq}}), \qquad
  c_i\T x \ge b_i \quad (i > m_{\mathrm{eq}}), \label{eq:pfeas}\\
  \lambda_i &\ge 0 \quad (i > m_{\mathrm{eq}}), \label{eq:dfeas}\\
  \lambda_i\,(c_i\T x - b_i) &= 0 \quad (i > m_{\mathrm{eq}}). \label{eq:comp}
\end{align}
Multipliers of equality constraints are unrestricted in sign, which is the one
asymmetry the algorithm has to carry through every subsequent formula.

\begin{proposition}[Sufficiency]\label{prop:sufficient}
If $(x,\lambda)$ satisfies \eqref{eq:stat}--\eqref{eq:comp} then $x$ is the
unique solution of~\eqref{eq:qp2}.
\end{proposition}

\begin{proof}
Let $\tilde x \in \Omega$. By strict convexity,
$f(\tilde x) \ge f(x) + (G x - a)\T(\tilde x - x)$, with equality only if
$\tilde x = x$. Substituting~\eqref{eq:stat},
\[
  (G x - a)\T(\tilde x - x)
  = \lambda\T\bigl(C\T \tilde x - C\T x\bigr)
  = \sum_i \lambda_i\bigl[(c_i\T\tilde x - b_i) - (c_i\T x - b_i)\bigr]
  = \sum_i \lambda_i (c_i\T \tilde x - b_i),
\]
the last step by~\eqref{eq:comp} for the inequalities and by $c_i\T x = b_i$ for
the equalities. Each surviving term is non-negative: for $i \le m_{\mathrm{eq}}$
the bracket vanishes, and for $i > m_{\mathrm{eq}}$ both factors are
non-negative by~\eqref{eq:pfeas} and~\eqref{eq:dfeas}. Hence
$f(\tilde x) \ge f(x)$ with equality only at $\tilde x = x$.
\end{proof}

Proposition~\ref{prop:sufficient} is the licence for everything in
Sections~\ref{sec:pdas} and~\ref{sec:sweep}. A point equipped with multipliers
passing the four conditions is not merely a plausible answer to be checked
further --- it is \emph{the} answer, and a method that produces candidates may
therefore be as unprincipled as one likes provided each candidate is tested. The
sufficiency is what makes an unguaranteed method safe rather than merely
usually right.

\subsection{The working-set subproblem}

Fix a working set $\Active$ of size $k$ with normals $A = C_{\cdot\Active}$ and
suppose $A$ has full column rank $k$. The \emph{working-set subproblem} is
\begin{equation}
  \min_{x} \; \tfrac12 x\T G x - a\T x
  \qquad \text{subject to} \qquad A\T x = b_\Active,
  \label{eq:sub}
\end{equation}
with every constraint in $\Active$ held as an equality regardless of its type in
the original program, and every constraint outside $\Active$ ignored.

\begin{proposition}[Solution of the subproblem]\label{prop:sub}
Problem~\eqref{eq:sub} has the unique solution and multipliers
\begin{equation}
  u = \bigl(A\T G^{-1} A\bigr)^{-1}\bigl(b_\Active - A\T x_u\bigr),
  \qquad
  x = x_u + G^{-1} A\, u,
  \qquad x_u := G^{-1} a .
  \label{eq:subsol}
\end{equation}
\end{proposition}

\begin{proof}
Stationarity for~\eqref{eq:sub} reads $Gx - a = Au$, so
$x = G^{-1}a + G^{-1}Au = x_u + G^{-1}Au$. Substituting into $A\T x = b_\Active$
gives $A\T x_u + (A\T G^{-1}A)u = b_\Active$. The matrix $A\T G^{-1}A$ is
positive definite because $G^{-1} \succ 0$ and $A$ has full column rank, hence
invertible, which yields~\eqref{eq:subsol}; the point so constructed is the unique
minimiser by strict convexity of $f$ on the affine
feasible set of~\eqref{eq:sub}.
\end{proof}

We call $A\T G^{-1} A$ the \emph{dual Hessian} of the working set. It appears
under three different names in what follows: as the matrix whose Cholesky factor
the solver carries (Section~\ref{sec:factorisation}), as the coefficient matrix of
the fast path's linear system (Section~\ref{sec:pdas}), and as the Gram matrix of
the least-squares problem the dual direction solves
(Lemma~\ref{lem:directions}). Its invertibility is equivalent to $A$
having full column rank, and that single fact is where every rank condition in
this paper comes from.

The pair $(x,u)$ of~\eqref{eq:subsol}, extended by zeros off the working set,
satisfies stationarity~\eqref{eq:stat} and complementarity~\eqref{eq:comp} by
construction --- the working-set constraints have zero slack and the rest have
zero multiplier. What it need not satisfy is the sign
condition~\eqref{eq:dfeas} on $\Active$'s inequalities, or primal
feasibility~\eqref{eq:pfeas} off $\Active$. The dual method maintains the first
and works towards the second.

\begin{remark}[Cold start]\label{rem:cold}
For $\Active = \emptyset$ the subproblem is unconstrained and~\eqref{eq:subsol}
degenerates to $x = x_u = G^{-1}a$ with an empty multiplier vector. The sign
condition holds vacuously, so this state satisfies the method's invariant with no
work beyond one Cholesky factorisation. That is the whole of the phase-one problem
for a dual method. The objective there is
$f(x_u) = \tfrac12 x_u\T G x_u - a\T x_u = -\tfrac12 a\T x_u$, which the
implementation uses as the initial value of the running objective.
\end{remark}

\subsection{The dual problem}

The method's name is earned by the following, which also explains why the
objective value can be tracked as a running total and why it increases.

\begin{proposition}[Dual function]\label{prop:dual}
The Lagrangian dual of~\eqref{eq:qp2} is the concave program
$\max\,\psi(\lambda)$ over $\lambda_i \ge 0$ $(i > m_{\mathrm{eq}})$, where
\begin{equation}
  \psi(\lambda) = -\tfrac12 (a + C\lambda)\T G^{-1} (a + C\lambda) + b\T\lambda .
  \label{eq:dualfn}
\end{equation}
If $(x,\lambda)$ satisfies stationarity~\eqref{eq:stat} and
complementarity~\eqref{eq:comp}, then $\psi(\lambda) = f(x)$.
\end{proposition}

\begin{proof}
Minimising $L(\cdot,\lambda)$ gives $x(\lambda) = G^{-1}(a + C\lambda)$ and,
writing $s = a + C\lambda$, $\psi(\lambda) = \tfrac12 s\T G^{-1}s - s\T G^{-1}s +
b\T\lambda$, which is~\eqref{eq:dualfn}. For the second claim \eqref{eq:stat}
gives $s = Gx$, and \eqref{eq:comp} with $c_i\T x = b_i$ on the equalities gives
$b\T\lambda = (Gx - a)\T x$; adding the two yields $f(x)$.
\end{proof}

So along any sequence of states of the form~\eqref{eq:subsol} the tracked value is
at once the primal objective at the subproblem minimiser and the dual objective at
the current multipliers. That is what makes the strict increase of
Proposition~\ref{prop:increment} progress towards optimality rather than
bookkeeping, and it gives the sandwich
$\psi(\lambda) \le f(x^\star) \le f(\tilde x)$: the method approaches the optimal
value from below, the mirror image of a primal active-set method.
